\title{Large Language Models Can Automatically Engineer Features for Few-Shot Tabular Learning}

\begin{abstract}
Large Language Models (LLMs) have demonstrated exceptional prowess in addressing complex, previously unencountered reasoning challenges, positioning them as highly promising tools for advancing tabular learning, a cornerstone of numerous practical domains. This work introduces a new in-context learning approach, \textsf{FeatLLM}, which leverages LLMs for automated feature engineering, thus generating data representations optimized for tabular prediction tasks. The engineered features serve as input for straightforward downstream models, such as linear regression, enabling highly effective few-shot learning. Distinctively, the \textsf{FeatLLM} strategy relies solely on the simple predictive model operating with the generated features during inference, eliminating recurrent interaction with the LLM at inference time—a common drawback in prior LLM-based approaches. Furthermore, \textsf{FeatLLM} only necessitates API-level access to LLMs, effectively sidestepping limitations imposed by prompt length. Experimental evaluations conducted across a diverse array of tabular datasets demonstrate that \textsf{FeatLLM} crafts high-quality, interpretable rules and, on average, outperforms comparable frameworks like TabLLM and STUNT by a significant margin (approximately 10\% improvement).
\end{abstract}

\section{Introduction}
Recently, Large Language Models (LLMs) have excelled at producing coherent and contextually appropriate outputs for novel tasks without requiring tailored fine-tuning~\cite{creswell2022selection,imani2023mathprompter,taylor2022galactica}. Through training on massive volumes of text, LLMs acquire extensive background knowledge, often obviating the dependence on expert input in various settings~\cite{Genomes2021,singhal2023large,wilcox2003role}, thereby enabling the automation of a broad spectrum of tasks, such as dialogue assessment~\cite{finch2023leveraging} and code correction~\cite{fan2023automated}. Notably, LLMs can interpret task prompts containing a brief description and a handful of examples, thus discerning context, identifying salient features, and focusing attention appropriately~\cite{brown2020language,zhao2023survey}. Such characteristics indicate that LLMs are well-equipped for performing advanced data analysis and, by extension, acting as sophisticated feature engineers.

Applying the rich, pre-trained knowledge and reasoning ability of LLMs to tabular data problems has garnered growing interest. For example, domain knowledge—like understanding the correlation between age and susceptibility to specific diseases—can enhance tasks such as disease prognosis. Current methodologies construct task definitions and feature explanations in natural language, serialize corresponding data, and provide this information as input to LLMs~\cite{dinh2022lift,wang2023anypredict}. Subsequent techniques may involve in-context learning~\cite{nam2023semi} or employ lightweight fine-tuning strategies~\cite{dinh2022lift,hegselmann2023tabllm}. Such utilization of LLMs' prior knowledge has consistently yielded better results than traditional tabular learning models, especially in settings with limited labeled data~\cite{hegselmann2023tabllm}.

Existing LLM-based approaches to tabular learning encounter notable drawbacks. Achieving end-to-end prediction necessitates at least a single LLM inference per data point, which introduces significant computational overhead. Moreover, attaining high predictive performance typically hinges on fine-tuning the underlying LLM—a requirement that is impractical for those models whose training infrastructure or tuning interfaces are restricted, especially in the case of recently introduced high-performing LLMs that are accessible only through limited inference APIs~\cite{anil2023palm,openai2023gpt4}. Another prevalent issue is the incompatibility with long prompts~\cite{liu2023lost}; when the tabular dataset contains numerous features such that the converted text input becomes extensive, these methods frequently become unworkable or their effectiveness diminishes. Collectively, these obstacles fundamentally impede the practical adoption of LLM-based models in real-world tabular learning scenarios.

To address these issues, our approach reconceptualizes the use of LLMs, assigning them the primary responsibility of feature engineering rather than direct end-to-end prediction. Specifically, we propose a novel in-context learning strategy, \textsf{FeatLLM}, which centers on extracting the underlying “criteria” that inform LLM decision-making. For example, when the task involves predicting a disease, the LLM can be leveraged to identify and articulate explicit rules—linking feature thresholds or patterns to the disease state. These derived criteria serve as the basis for new features, effectively substituting original variables in downstream tasks. This strategy transitions the LLM’s primary contribution to feature creation, permitting the use of simpler downstream predictive models and resulting in improved inference efficiency. The capacity of LLMs for in-context learning enables us to generate new features purely through the inclusion of illustrative examples within the prompt, eliminating the need for explicit model training. Furthermore, integrating ensemble techniques such as feature bagging~\cite{breiman2001random} helps control prompt length, thus mitigating issues arising from large input sizes.

\textsf{FeatLLM} decomposes the feature engineering prompt into two distinct stages: (1) analyzing the problem context to infer relationships between features and target variables; and (2) employing insights from the initial analysis, along with a limited number of examples, to construct explicit rules for class differentiation. This sequential reasoning methodology encourages the LLMs to disregard irrelevant features in the process of rule formulation. The resulting rules are implemented on the data—interpreted as code—and each rule yields a binary feature flagging whether the data point satisfies the respective condition. Subsequently, a simple, low-complexity predictive model is trained to estimate class probabilities using these constructed features. To enhance robustness, an ensemble strategy incorporating repeated feature generation and feature bagging is used. This approach also addresses issues stemming from oversized prompts by tuning the selection of features and data samples during bagging. Since \textsf{FeatLLM} operates without the need for further model training and has low inference requirements, leveraging LLMs in diverse practical scenarios is achievable.

Our empirical study assesses \textsf{FeatLLM} on 13 tabular datasets under low-shot learning conditions, demonstrating both its robustness and efficacy. The results indicate that LLMs successfully integrate domain expertise with insights derived from the available data to produce valuable features. Across a wide range of scenarios, our method surpasses state-of-the-art few-shot learning baselines. The source code is publicly available via an anonymized GitHub repository at~\texttt{https://github.com/Sungwon-Han/FeatLLM}.

\section{Related Work}

\subsection{Few-Shot Learning with Tabular Data} Few-shot learning has made it possible to obtain transferable and robust data representations while significantly reducing labeling expenses~\cite{chen2018closer}. Although early efforts centered on images, the few-shot learning paradigm has shown its effectiveness across multiple modalities such as text, audio, and more recently, tabular data~\cite{majumder2022few,nam2023stunt,schick2021exploiting}. Working with only a handful of labeled samples, however, increases susceptibility to overfitting or detecting misleading patterns~\cite{bartlett2021deep}. To mitigate this, contemporary research leverages supplementary information sources and integrates prior domain knowledge, thus instilling suitable inductive biases during model development. For example, semi-supervised approaches utilize abundant unlabeled tabular data in conjunction with thoughtful augmentation~\cite{ucar2021subtab,yoon2020vime}, while some strategies seek advantageous initial conditions for models through broad, task-independent unsupervised pre-training~\cite{somepalli2021saint}. In these cases, unsupervised learning techniques include masking or perturbing parts of the input and employing reconstruction-based losses~\cite{arik2021tabnet,majmundar2022met}, or relying on data augmentation methods to implement contrastive learning frameworks~\cite{bahri2022scarf,chen2020simple}. Nevertheless, the structural diversity inherent to tabular data complicates the design of generic augmentations, resulting in limited efficacy of these methods in few-shot regimes~\cite{nam2023stunt}. 

Recent advances advocate for training models across an array of heterogeneous tabular datasets, which may not necessarily be related to the final application domain, and subsequently utilizing transfer learning for specific target tasks~\cite{zhang2023generative,levin2022transfer,zhu2023xtab}. As an example, TransTab~\cite{wang2022transtab} utilizes transformer architectures to learn semantic connections among columns by considering column names, rather than relying solely on table entries. Insights derived from these various columns and tables have been shown to benefit zero-shot and few-shot transfer to new tabular learning problems. Similarly, TabPFN~\cite{hollmann2022tabpfn} generates synthetic datasets by blending multiple distribution types and employs these for pre-training a Bayesian neural network. During inference, both training and test instances from the new task are provided simultaneously to the pre-trained model without the need for additional fine-tuning. Prior knowledge acquired across such diverse datasets enables superior performance to conventional tree-based methods, and this transfer has proven particularly effective in low-shot learning contexts.

\subsection{Language-Interfaced Tabular Learning} The utilization of large language models (LLMs)~\cite{anil2023palm,openai2023gpt4} presents an alternative strategy for integrating external knowledge. Owing to pre-training on expansive and varied textual datasets, LLMs are adept at generalizing to new, unseen tasks and have shown impressive results across many areas~\cite{brown2020language}. Recent work explores methods for exploiting the extensive latent knowledge within LLMs in the context of tabular data modeling. A common practice is to convert tabular data into textual formats and present them as input prompts for LLMs~\cite{dinh2022lift,hegselmann2023tabllm,wang2023anypredict}. By embedding the table values, feature annotations, and the problem descriptions within the prompt, these systems infer dependencies between features and the task to accomplish prediction or classification objectives. Progress in this direction includes model adaptation techniques like parameter-efficient fine-tuning with LoRA~\cite{hu2021lora} or IA3~\cite{liu2022few}, as well as leveraging in-context learning where demonstration examples are included directly in the prompts, enabling inference without explicit model re-training~\cite{dinh2022lift,hegselmann2023tabllm,nam2023semi,slack2023tablet}.

Although the previously mentioned approaches have demonstrated efficacy in enhancing few-shot tabular learning, their reliance on routing inference exclusively through the LLM introduces obstacles for deployment in industrial contexts. This work seeks to move away from the traditional black-box paradigm, in which every individual sample’s prediction is derived end-to-end using an LLM. Instead, the approach pivots to utilizing the LLM for directly inferring class-specific conditions or rules. \textsf{FeatLLM} translates the identified rules into programmatic code to generate novel features, thereby facilitating predictions independent of further LLM calls. This leverages in-context learning to distill rules from existing background knowledge in conjunction with the provided few-shot examples.

\section{Methods}
\textsf{FeatLLM} is engineered to uncover “rules”—that is, the attribute conditions linked to each output class—from the small set of provided examples, as illustrated in Figure~\ref{fig:main_model}. These rules are derived with the aid of an LLM and are subsequently converted into programmatic conditions to form new binary-valued features for each sample, capturing whether each rule is fulfilled. The resulting binary feature set for a sample is then processed via a dot product with non-negative, trainable weights, yielding class probability estimates. Training this model allows for the estimation of the relevance of each extracted feature in determining class probabilities (see Section~\ref{Sec:training} for details). This workflow is executed repeatedly within a bagging framework, and outputs are integrated through ensembling. The problem definition and details for each component are outlined below.

\vspace*{-0.12in}
\paragraph{Problem formulation.} Consider a tabular dataset comprising $N$ labeled data points, denoted as $\mathcal{D} = \{(\mathbf{x}^i, \mathbf{y}^i)\}_{i=1}^{N}$. Each sample $\mathbf{x}^i$ is represented as a $d$-dimensional input feature vector, and the dataset provides human-readable feature names $F = \{f_j\}_{j=1}^{d}$, with accompanying definitions available separately. Within the classification framework, each label $\mathbf{y}^i$ belongs to a set of predefined categories. In $k$-shot learning scenarios, the training set consists of only $k$ randomly selected labeled instances, with $k < N$.

\subsection{Prompt Design for Extracting Rules}
\label{Sec:prompt_design}
To improve the quality of rule extraction using an LLM, we shape the reasoning trajectory to reflect expert approaches to tabular data analysis. The prompt we develop incorporates three fundamental elements (see Fig.~\ref{fig:prompt_for_rule} and Appendix~\ref{sec:example_rule}):

\vspace*{-0.12in}
\paragraph{Basic information description.} This section supplies key context required for task completion (highlighted in orange in Fig.~\ref{fig:prompt_for_rule}). It consists of a specification of the problem, including label information, feature descriptions, and illustrative examples. The task prompt is presented in question form (e.g., ``Does this patient's myocardial infarction complications data indicate chronic heart failure? Yes or no?''; further examples are available in Appendix~\ref{sec:task_desc}). Each feature is described along with its data type and, if appropriate, its definition; categorical variables may be supplemented with example values. For illustrative purposes, a handful of training examples and their corresponding labels are provided as text sequences, formatted in alignment with established practices in prior research~\cite{dinh2022lift,hegselmann2023tabllm}:

\begin{align}
&\text{Serialize}(\mathbf{x}^i,\mathbf{y}^i, F) = \nonumber \\ 
&``\ f_1\ \text{is}\ \mathbf{x}^i_1.\ ... \ f_d\  \text{is}\  \mathbf{x}^i_d.\ \ \text{Answer:}\  \mathbf{y}^i\ ”
\end{align}

\vspace*{-0.12in}
\paragraph{Reasoning instruction.} To promote deeper reasoning abilities within the LLM, we incorporate explicit reasoning instructions (see blue text in Fig.~\ref{fig:prompt_for_rule}). The reasoning instruction opens with a sentence modeled after the chain-of-thought paradigm~\cite{wei2022chain}, after which we structure the inference procedure into two distinct phases. In the initial phase, the LLM is tasked with deducing potential causal relationships or tendencies linking features to the given task description, deliberately excluding consideration of example demonstrations. This encourages reliance on broad, commonsense knowledge, thereby maximizing the LLM’s use of background understanding relevant to the task. The second phase involves utilizing both the insights gained in the first phase and example demonstrations to infer class-specific rules. Adopting this bifurcated strategy is intended to steer the model away from learning spurious associations related to irrelevant data fields and encourages focus on features of higher importance. To avoid both underfitting from too few rules and unnecessarily long prompts resulting from too many, we impose a predetermined limit on the number of rules generated for each class (specifically, 10)\footnote{Section~\ref{sec:analysis} contains further discussion regarding the selection of rule quantity.}. Furthermore, the rule-generation process is informed by the LLM’s awareness of each feature's datatype as specified in the fundamental description, ensuring semantic coherence and type consistency.

\vspace*{-0.12in}
\paragraph{Response instruction.} To support effective extraction and downstream utilization of the derived rules, we direct the LLM through tailored response instructions (see yellow text in Fig.~\ref{fig:prompt_for_rule}). These directions specify the format in which responses for each target class should be organized (i.e., response structure), as well as the prescribed template for rule articulation (i.e., rule structure). Such explicit guidance enables the LLM to choose formats appropriate to each feature type. In the absence of these guidelines, the model is otherwise permitted to integrate several rules via logical connectives, such as AND or OR. An illustration of the resultant output can be found in Appendix~\ref{sec:outcome-rule}.

\subsection{Inferring Class Likelihood via Rules}
\label{Sec:training}

\paragraph{Parsing rules for feature generation.} In this stage, we transform the rules derived earlier into new binary features. For every class, these binary indicators specify whether a sample fulfills the class-specific rules (taking values of '0' or '1'). Such features provide information about a sample's membership likelihood in each respective class. Nevertheless, since these rules are originally articulated in natural language by the LLM, a dedicated parsing method is necessary for automatic feature extraction. Although we enforce formatting constraints on both the responses and rules to aid in systematic parsing, LLM outputs can still exhibit unexpected or inconsistent syntactic modifications~\cite{kovavc2023large}. For example, instead of symbols like ``$>$" or ``$<$", the LLM may use words such as ``greater than" or ``less than", thereby complicating the parsing task. 

To overcome the parsing difficulties posed by noisy text, rather than implementing intricate parsing algorithms, we make use of the LLM itself for this step. The LLM demonstrates strong semantic understanding regardless of surface-level textual variations, and it can reliably produce straightforward code when given structured instructions (thanks to its training on code datasets). Accordingly, our prompts specify the function name, descriptions of input and output formats, and the relevant extracted rules, all of which are supplied to the LLM. Figures~\ref{fig:prompt_for_parsing} and~\ref{sec:example_parsing}-\ref{sec:example_parse_outcome} in the Appendix illustrate sample prompts and resulting outputs. The Python \texttt{exec()} function is then used to execute the generated code for data processing.

\vspace*{-0.12in}
\paragraph{Inferring class likelihood.} The newly constructed binary rule-based features allow for a straightforward estimation of class likelihood by tallying the number of satisfied rules per class (that is, summing the binary indicators for each class). Still, as different features and rules may contribute unevenly to prediction, their weights must be inferred from training data. To achieve this, we employ standard machine learning (ML) algorithms to estimate the relative importance of each binary feature for every class. For instance, a linear model without an intercept can be adopted. Formally, for sample $\mathbf{x}^i$, let $\mathbf{z}_k^i \in \{0, 1\}^R$ denote its vector of binary features for class $k$, where $R$ is the rule count per class, and $c$ is the number of classes. We compute the probability for each class $\mathbf{p}^i$ as follows:

\begin{align}
\text{logit}_k^i &= \max(\mathbf{w}_k, 0) \cdot \mathbf{z}_k^i, \nonumber \\
\mathbf{p}^i &= \text{Softmax}({[\text{logit}_{1}^i, ..., \text{logit}_{c}^i]}). \label{eq:infer_prob}
\end{align}

In this context, $\mathbf{w}_k$ denote the adjustable parameters within the linear model, which quantify the relative contribution of each feature. By constraining these parameters to a positive domain, we guarantee that the LLM-derived features do not detract from the probability of assigning an instance to a given class during model training. Following this, the resulting logit vector is transformed into a probability distribution over classes by applying the softmax function. The weights $\mathbf{w}_k$ are optimized by minimizing the cross-entropy loss, employing a limited set of labeled examples in a few-shot learning regime.

\vspace*{-0.12in}
\paragraph{Ensembling with bagging.} To further enhance predictive performance, we iterate the entire pipeline multiple times, thus generating several independent inference models, and aggregate their outputs through an ensemble mechanism. Specifically, for $T$ independent runs producing sets of learned weights $\{\mathbf{w}[t]\}_{t=1}^T$, the final class prediction is based on the average class probabilities $\mathbf{p}^i$ evaluated for each $\mathbf{w}[t]$, consistent with Eq.~\ref{eq:infer_prob}.

To foster diversity among the ensemble's constituent models, randomness is injected at two levels. First, the LLM inference utilizes a relatively elevated temperature parameter, promoting variation in rule synthesis. Second, the sequence of few-shot examples is randomized for each trial, leveraging findings that example ordering can influence LLM outputs~\cite{lu2022fantastically}\footnote{Further discussion of rule diversity is available in Appendix~\ref{sec:diversity}}. Additionally, to maximize heterogeneity in prediction and bolster robustness, we incorporate bagging by selecting distinct feature or instance subsets in each run—a strategy that proves especially beneficial as the number of available training samples or features increases. This ensemble strategy brings two chief benefits: (1) If certain LLM-generated rule sets are flawed in some trials, accurate rule sets from others can offset them, thereby bolstering the overall reliability of predictions. This property is linked to the notion of LLM self-consistency~\cite{wang2022self}, since rules consistently generated across runs tend to be more reliable. (2) The approach helps to mitigate the limitations imposed by the LLM's input length constraints; in situations where tabular data are too voluminous for a single prompt, bagging reduces the prompt size required per run, maintaining strong model performance. Rather than relying on a single, possibly incomplete rule set, this ensemble of multiple, comparatively weak learners collectively compensates for gaps in feature or instance coverage present in individual model runs.

\section{Experiments}
We assess the performance of \textsf{FeatLLM} across several tabular datasets and carry out a range of studies to better understand the operation of our proposed approach. For brevity, supplementary experiments—such as those varying the choice of LLMs, performing qualitative assessments, and additional analyses—are documented in the Appendix.

\subsection{Performance Evaluation}
\paragraph{Datasets.}
Our experimental setup encompasses 13 tabular datasets designed for either binary or multi-class classification tasks: (1) Adult~\cite{asuncion2007uci} targets the prediction of whether a person’s income exceeds \$50,000 annually; (2) Bank~\cite{moro2014data} predicts a customer’s likelihood to subscribe to a term deposit; (3) Blood~\cite{yeh2009knowledge} anticipates whether blood donors will return for another donation; (4) Car~\cite{kadra2021well} is used for determining automobile quality; (5) Communities~\cite{misc_communities_and_crime_183} predicts the crime rate within a region; (6) Credit-g~\cite{kadra2021well} is aimed at classifying individuals as good or bad credit risks; (7) Diabetes\footnote{\texttt{kaggle.com/datasets/uciml/pima-indians-diabetes-database}} identifies diabetic status in individuals; (8) Heart\footnote{\texttt{kaggle.com/datasets/fedesoriano/heart-failure-prediction}} forecasts the likelihood of coronary artery disease; (9) Myocardial~\cite{misc_myocardial_infarction_complications_579} ascertains the risk of chronic heart failure in patients.

In addition, we include both contemporary and artificially-generated datasets that are seldom featured in literature and excluded from LLM pre-training: (10) Cultivars estimates the yield potential of different soybean cultivars\footnote{\texttt{archive.ics.uci.edu/dataset/913}}; (11) NHANES seeks to classify individuals’ age groups based on provided records\footnote{\texttt{archive.ics.uci.edu/dataset/887}}; (12) Sequence-type categorizes input sequences as either arithmetic, geometric, Fibonacci, or Collatz types; (13) Solution-mix involves predicting whether the mixture of four solutions surpasses a 0.5 percentage concentration threshold. The first two datasets were recently made available on the UCI Machine Learning Repository post-September 2023, confirming their absence from the pre-training data of the LLM API (gpt-3.5-turbo-0613) leveraged in this work. The last two are synthetically constructed, guaranteeing that their contents were unknown to the LLM at evaluation time.

A range of scales and complexities are represented in these datasets, with more information presented in Table~\ref{Tab:basic-desc}. Each includes clear attribute labels and descriptions. Certain datasets—such as `Communities' and `Myocardial'—comprise over 100 features, presenting notable difficulty for LLMs due to limited context window capacity.

\vspace*{-0.12in}
\paragraph{Baselines.}
We benchmark our approach against nine baseline methods. The initial three are widely used algorithms for tabular data: (1) Logistic regression (LogReg), (2) XGBoost~\cite{chen2016xgboost}, and (3) RandomForest~\cite{ho1995random}. Following these, two baselines assume access to abundant unlabeled data: (4) SCARF~\cite{bahri2022scarf} and (5) STUNT~\cite{nam2023stunt}. However, collecting large-scale unlabeled data may not always be practical in many scenarios. Another recent baseline, (6) TabPFN~\cite{hollmann2022tabpfn}, uses a variety of synthetic data distributions to pretrain the model. Lastly, three baselines utilize LLM-based strategies: (7) In-context learning (In-context)~\cite{wei2022emergent}, (8) TABLET~\cite{slack2023tablet}, and (9) TabLLM~\cite{hegselmann2023tabllm}. In-context learning operates by incorporating few-shot examples directly into the input prompt, without model adaptation. TABLET supplements the prompt with external rule-based and prototypical hints derived from a separate classifier to bolster inference. TabLLM uses parameter-efficient tuning, specifically the IA3 technique~\cite{liu2022few}, to adapt the LLM using a limited set of training instances.

\vspace*{-0.12in}
\paragraph{Implementation details.}
\textsf{FeatLLM} employs GPT-3.5 as its underlying LLM, but its framework remains flexible and compatible with various LLM backbones (additional results with alternative backbones can be found in Appendix~\ref{sec:backbones}). During LLM inference, the temperature parameter is maintained at 0.5, while the top-p parameter follows the API's default value of 1. For rule extraction, the number of ensembles and rules is set to 20 and 10, respectively. The influence of these hyper-parameters is illustrated in Figure~\ref{fig:hyperparameter2} and further discussed in Appendix~\ref{sec:hyper-parameter}. The linear model is optimized using Adam, adopting a learning rate of 0.01, and trained over 200 epochs. Selection of the best epoch is achieved via $k$-fold cross-validation. When re-implementing baseline methods, GPT-3.5 is used for in-context learning setups, while T0~\cite{sanh2022multitask} is used with TabLLM to be consistent with its original configuration. Notably, TabLLM only permits LLMs for which public checkpoints exist, which excludes GPT-3.5 from its options. As for conventional machine learning models such as LogReg, XGBoost, and RandomForest, we perform parameter selection through Grid-search combined with $k$-fold cross-validation. The parameter $k$ is chosen as either 2 or 4, ensuring every class is represented in the training splits. Appendix~\ref{sec:baselines} provides additional implementation specifics.

\vspace*{-0.12in}
\paragraph{Results.}
Tables~\ref{tab:main1} and \ref{tab:main2} display a comparative analysis of performance metrics across multiple datasets. Each experiment was conducted three times with distinct random splits for training and testing, reporting the mean and standard deviation of AUC outcomes. \textsf{FeatLLM} reliably emerges as either the best or second-best model relative to all baseline competitors. Notably, our method, utilizing only 4 training examples, matches the performance of standard tabular baselines that require 32 examples, as illustrated in Fig.~\ref{fig:compares}. Although models such as STUNT and TabPFN achieved marked gains on some individual datasets, their aggregate improvement over traditional machine learning approaches remains modest. Furthermore, LLM-based solutions—like in-context learning, TABLET, or TabLLM—struggled with inference tasks on datasets characterized by high feature dimensionality. By incorporating both bagging and ensemble methods, our framework consistently delivered robust results even in high-dimensional scenarios. It is also worth highlighting that our ensemble-based strategy could be integrated into preexisting LLM-based pipelines. However, such integration would effectively double inference operations per data sample, dramatically escalating computational demands and making it less feasible in practice.

\subsection{Analysis \& Discussion}
\label{sec:analysis}
\paragraph{Ablation study.} We assess the influence of the main components on model effectiveness via a series of ablation experiments. These consider: (1) \textsf{-Tuning}: excluding the weight tuning in the linear model, thereby aggregating binary features to generate logits directly; (2) \textsf{-Ensemble}: removing the ensemble process; (3) \textsf{-Description}: leaving out feature descriptions; and (4) \textsf{-Reasoning}: omitting Step-1 of the reasoning instruction used during rule generation.

The summary of AUC variations, averaged over all datasets and caused by these ablations, is provided in Table~\ref{tab:ablation}. In every scenario, omission of a specific component consistently causes performance declines. Notably, the advantage conferred by weight tuning amplifies as the shot count increases; this pattern implies that a larger dataset supports better rule importance estimation, heightening the positive impact of tuning. Conversely, the importance of detailed feature descriptions and explicit reasoning instructions is more pronounced when fewer training examples are available, underlining the significance of leveraging the LLM’s prior knowledge for performance boosts in low-data regimes. Finally, the proposed ensemble technique demonstrates steady improvements to overall model effectiveness across all experimental settings.

\vspace*{-0.12in}
\paragraph{Training \& inference time comparison.} We evaluate and contrast the training and inference durations across various methods. All experiments utilize the Adult dataset with a training set size of 16 shots. A single A100 GPU is generally employed for these benchmarks, except in the case of TabLLM, which leverages four A100 GPUs to enable model parallelism. Table~\ref{tab:cost} displays the corresponding runtime results. Importantly, \textsf{FeatLLM} achieves notably efficient inference times, on par with traditional baselines. By comparison, inference with other LLM-based techniques, such as In-context learning, TABLET, and TabLLM, requires considerably more time. As for training time, \textsf{FeatLLM} is similar to other LLM-based strategies—only 30 API queries are necessary, ensuring that the method remains cost-effective and the process can be easily parallelized if desired.

\vspace*{-0.12in}
\paragraph{Quality of features generated by \textsf{FeatLLM}.}
To assess the utility of the rule-derived features created by \textsf{FeatLLM}, we perform a straightforward experiment measuring their contribution to practical task performance. Specifically, for each dataset, we calculate the average AUROC between each newly introduced feature and the target variable, subsequently reporting the proportion of features that attain an AUROC exceeding 0.5. Table~\ref{tab:quality} compiles these findings. The results indicate that most generated rules capture pertinent information related to the predictive task (see the ``Before tuning'' column). Additionally, during linear model fine-tuning, features exhibiting minimal relevance—defined as those whose learned coefficients fall below a specified threshold—are systematically removed (see the ``After tuning'' column). The threshold is fixed at 0.1, based on the observed weight distribution.

\vspace*{-0.12in}
\paragraph{Effect of adding spurious correlations.} To evaluate how well different approaches eliminate noisy, spurious correlations under low-data conditions, we performed an assessment. Our experimental setup relied on the Heart dataset, augmented by the random inclusion of columns from the Adult dataset that have no relevance to the Heart dataset's classification objective. The analysis examined model performance as we systematically increased the number of extraneous columns. To ensure a fair evaluation, all methods operated under the constraint of only 8 labeled samples, with no access to any unlabeled data, which means techniques like SCARF or STUNT were not included in the comparison.

The impact of these spurious correlations on model performance is illustrated in Figure~\ref{fig:spurious}. Notably, \textsf{FeatLLM} exhibited the smallest decline in accuracy as irrelevant features increased, compared to the baselines. This robustness likely stems from our approach, which emphasizes the interplay between the target task and the chosen features using existing knowledge during rule derivation. As a result, rules seldom reference irrelevant columns, leading to more stable performance; indeed, we found that only about 1-2\% of extracted rules incorporated the noise variables. However, we also noticed that \textsf{FeatLLM} can still be misled by spurious correlations and conflate causal associations between tasks and features when unrelated but contextually similar columns are appended at random (refer to Appendix~\ref{sec:spurious-appendix} for further discussion).

\vspace*{-0.12in}
\paragraph{Hyper-parameter analysis}
\textsf{FeatLLM} incorporates two primary hyper-parameters: the ensemble count and the quantity of rules extracted per class during each LLM interaction. To understand their effects, we systematically varied these hyper-parameters and observed resulting performance changes. Our findings indicate that increasing the ensemble count initially yields improved outcomes; however, after reaching approximately 20 ensembles, performance plateaus (refer to Fig.~\ref{fig:appendix-hyperparameter1} in Appendix). Concerning rule count per class, there were no statistically significant performance variations, but a value of 10 was identified as optimal, balancing prompt size with effectiveness (see Fig.~\ref{fig:hyperparameter2}). We hypothesize that incorporating additional rules expands the input’s dimensionality, increasing informational content while simultaneously risking overfitting in low-shot scenarios.

\section{Conclusion}
We introduce \textsf{FeatLLM}, which sets a new benchmark in LLM-driven few-shot learning for tabular data. Unlike approaches that use LLMs merely as direct predictors for each record, \textsf{FeatLLM} leverages LLMs to devise predictive criteria in the role of feature engineers. With its rule generation strategy and ensemble approach via bagging, \textsf{FeatLLM} incurs minimal inference overhead, operates using only API-level access with no training required, and effectively sidesteps data feature dimensionality limitations. This method consistently yields strong predictive results, establishing itself as a highly efficient solution for applying LLMs to practical tabular datasets.

\textsf{FeatLLM} is currently tailored for scenarios characterized by limited data availability. Moving forward, our aim is to enhance its functionality to accommodate larger datasets, explore alternative forms of feature construction beyond rules, and further improve interpretability, thereby broadening its applicability across more diverse tasks.

\section{Broader Impact}
The proposed \textsf{FeatLLM} framework seeks to exploit LLMs’ inherent prior knowledge to achieve tabular learning outcomes surpassing those obtainable through traditional supervised methods. By augmenting learning with such prior knowledge, the approach significantly alleviates overfitting risks stemming from scant training data, thereby promoting more data-efficient learning. \textsf{FeatLLM} offers substantial benefits to practitioners in domains such as finance and healthcare, where labeled data are often scarce or costly to obtain, and LLMs’ extensive pretraining on related textual data can be invaluable. An essential direction for future research is to scrutinize the broader societal implications, particularly regarding the propagation of existing biases and misinformation embedded in the LLMs’ prior knowledge, as this information could potentially influence or override insights drawn from observed labeled samples.

\end{document}